\documentclass[12pt, letterpaper]{article}

%% Build from the paper/ directory:
%%   cd paper && pdflatex main.tex
%% or use the root build system:
%%   make paper          (from repo root)
%%   ./build.sh paper    (POSIX)
%%   build.cmd paper     (Windows cmd)

% packages.tex
\usepackage{amsmath, amssymb, amsthm}
\usepackage{graphicx}
\usepackage{xcolor}
\usepackage[
    hidelinks,
    pdftitle={Optical-Axis Birefringence on the A=1 Discrete Causal Lattice},
    pdfauthor={Jack Menendez},
    pdfsubject={A=1 Discrete Causal Lattice -- optical-axis birefringence (Paper IV)},
    pdfkeywords={A=1 Discrete Causal Lattice, optical-axis birefringence, kinematic birefringence, gauge-sector birefringence, (1,1,-1) optical axis, Dirac structure}
]{hyperref}
\usepackage{booktabs}
\usepackage{array}
\usepackage{geometry}
\usepackage{adjustbox}
\usepackage{longtable}
\usepackage{setspace}
\usepackage{caption}
\usepackage{subcaption}
\usepackage{float}
\onehalfspacing

\input{macros/commands}

%% Page geometry (also set here as a fallback in case the package path fails)
\geometry{top=1in, bottom=1in, left=1in, right=1in}

%% ---------------------------------------------------------------
%% TITLE PAGE
%% ---------------------------------------------------------------
%% TODO: replace title, subtitle, author, ORCID, email.
%% TODO: at release time, fill in the \thanks{} block with version,
%%       Zenodo DOI, GitHub URL, and any series identifier.

\title{Optical-Axis Anisotropy on the A=1 Discrete Causal Lattice:\\
       Birefringence Cancellation and a Single-Domain No-Go%
       \thanks{Version 1.0.  Paper~IV of the A=1 Discrete Causal
       Lattice series; takes up follow-on item~\#14
       (\emph{Balanced Equations and Birefringent Channels}) of
       Paper~I's \texttt{notes/follow\_on\_implications.md} catalogue. This
       paper subsumes the former companion study of the electric induced-action
       block (Paper~VIII), merged here since the shared optical axis is the
       natural unity of the two effects.
       Builds on \emph{Geometry First}
       (Paper~I,
       \href{https://doi.org/10.5281/zenodo.20613677}{DOI:
       10.5281/zenodo.20613677}) and \emph{Geometry Forces Physics}
       (Paper~II,
       \href{https://doi.org/10.5281/zenodo.20292158}{DOI:
       10.5281/zenodo.20292158}).
       Paper~III (tidal ionization) is deferred pending the
       $\delta p_\text{min}$ investigation.
       \href{https://github.com/JackDMenendez/dcl-paper-04-optical-axis-birefringence}%
            {GitHub: JackDMenendez/dcl-paper-04-optical-axis-birefringence}.
       \href{https://doi.org/10.5281/zenodo.21435952}{DOI:
            10.5281/zenodo.21435952}.}}
\author{Jack Menendez\\
        \small Eugene, OR\\
        \small ORCID: \href{https://orcid.org/0009-0003-1166-307X}{0009-0003-1166-307X}\\
        \small \href{mailto:jack@geometryinducedphysics.org}{jack@geometryinducedphysics.org}\\
        \small Web: \href{https://geometryinducedphysics.org/}{geometryinducedphysics.org}}
\date{\today}

\begin{document}

\maketitle

%% ---------------------------------------------------------------
%% FRONT-MATTER NOTICES
%% ---------------------------------------------------------------
\noindent\textbf{License.}  Paper text and figures: Creative Commons
Attribution 4.0 (CC BY 4.0).  Source code in the companion
repository: MIT License.

\medskip
\noindent\textbf{Data availability.}  The LaTeX source for this paper
is publicly available at
\href{https://github.com/JackDMenendez/dcl-paper-04-optical-axis-birefringence}%
     {\nolinkurl{github.com/JackDMenendez/dcl-paper-04-optical-axis-birefringence}}.
%% At release time, point to the tag and the Zenodo DOI:
%% The exact commit corresponding to this release is the vX.Y tag,
%% archived at
%% \href{https://doi.org/10.5281/zenodo.XXXXXXXX}%
%%      {\nolinkurl{doi:10.5281/zenodo.XXXXXXXX}}.

\medskip
\noindent\textbf{Keywords.}  A=1 Discrete Causal Lattice; optical-axis
birefringence; kinematic birefringence; gauge-sector birefringence;
$(1,1,-1)$ optical axis; Dirac structure; lattice anisotropy;
multi-channel concordance.

\newpage

\begin{abstract}
    % abstract.tex
The A=1 discrete causal lattice of Paper~I~\cite{menendez2026geometryfirst} is, by
construction, a birefringent medium: its bipartite hop geometry singles out the
optical axis $\hat{\mathbf{n}}_* = (1,1,-1)/\sqrt3$---the one cube body-diagonal
absent from the hop set---so propagation along the axis is inequivalent to
propagation across it. We work this structural fact to closed form in two sectors
that share that single axis. In the \emph{kinematic} sector, the single-tick spinor
propagator yields an exact ordinary/extraordinary speed split: the dispersion
depends only on the wavevector component perpendicular to $\hat{\mathbf{n}}_*$, with
effective-mass tensor $M_\text{eff} = \tfrac43 P_\perp$, a group-velocity split
$\Delta v = \tfrac23\sin\omega\,|\mathbf{k}|$ in the massive sector, and an
$\omega$-independent dispersion-curvature split ($4/3$ versus $0$) carried by the
photon and massive sectors alike. This is the paper's falsifiable dispersion
prediction: fixed in direction by $\hat{\mathbf{n}}_*$ and parameter-free up to the
lattice scale. In the \emph{gauge} sector, integrating the matter tokens out of the
A=1 dynamics induces an anisotropic Maxwell action whose permittivity and
inverse-permeability tensors, $\boldsymbol{\varepsilon}$ (eigenvalues $\{1,4,4\}$)
and $\boldsymbol{\mu}^{-1}$ (eigenvalues $\{4,4,16\}$)---both uniaxial about
$\hat{\mathbf{n}}_*$, in opposite sense---are read directly off the running engine's
hop holonomies. Because $\boldsymbol{\mu}^{-1} = \operatorname{adj}
(\boldsymbol{\varepsilon})$, the photon dispersion is a double transverse root and
the gauge-sector vacuum \emph{birefringence cancels exactly} (a null we confirm
numerically to $10^{-15}$), leaving a factor-$\sim2$ common-mode speed anisotropy
whose observability we leave open. The gauge-sector derivation and the adjugate
cancellation theorem are the companion Paper~VIII~\cite{menendez2026electricblock};
here they are confirmed independently from the engine. Every claim's scope, evidence,
and status is recorded in the audit table (Table~\ref{tab:audit}).

\end{abstract}

%% ---------------------------------------------------------------
%% RELEASE NOTES (only at release time)
%% ---------------------------------------------------------------
%% \bigskip
%% \noindent\textbf{Release notes for vX.Y (since vX.Y-1).}
%% \begin{itemize}
%% \item Item 1.
%% \item Item 2.
%% \end{itemize}

%% \emph{Status legend used throughout.}  \texttt{PASS} / \texttt{PART}
%% / \texttt{STUB} / \texttt{FAIL} -- see Table~\ref{tab:audit}.

\newpage
\tableofcontents
\newpage

%% ============================================================
%% PART I: <name>
%% ============================================================
\section{Introduction}
% introduction.tex
% Orientation for Paper IV: the lattice as a birefringent medium with a
% geometry-fixed optical axis; the two channels (kinematic + gauge); the
% companion structure with Paper VIII; scope. The audit table is included in
% the Scope subsection so the reader meets it before any status-dependent claim.

\label{sec:introduction}

A crystal is birefringent when its atomic lattice singles out an axis, so that the
speed of light depends on direction and polarization. The A=1 discrete causal
lattice of Paper~I~\cite{menendez2026geometryfirst} is birefringent for the same
reason, but the ``crystal'' is spacetime itself. Its bipartite hop set---the three
RGB vectors $\mathbf{V}_1=(1,1,1)$, $\mathbf{V}_2=(1,-1,-1)$,
$\mathbf{V}_3=(-1,1,-1)$ and their CMY negatives---uses three of the four cube
body-diagonals and leaves the fourth, $\hat{\mathbf{n}}_* = (1,1,-1)/\sqrt3$, out.
That single omission breaks the cubic point group $O_h$ to the trigonal $D_{3d}$
about $\hat{\mathbf{n}}_*$ and fixes an optical axis geometrically, with no new
postulate. The optical-axis birefringence is one of the series' standing
predictions; this paper works it to closed form.

The effect appears in two sectors, and they share the \emph{same} axis. The
\emph{kinematic} channel (Sec.~\ref{sec:kinematic_channel}) follows from Paper~I's
single-tick spinor propagator: the entire wavevector dependence passes through one
scalar structure factor whose modulus is flat along $\hat{\mathbf{n}}_*$ and curved
perpendicular to it, giving a closed-form ordinary/extraordinary speed split. This
is the falsifiable, direction-fixed dispersion prediction the paper contributes, and
its observability---that the anisotropy survives the per-tick A=1 renormalization
into observable amplitudes---is established numerically. The \emph{gauge} channel
(Sec.~\ref{sec:gauge_sector}) is the response of the vacuum: integrating the matter
tokens out of the dynamics in a background electromagnetic field induces an
anisotropic Maxwell action whose permittivity $\boldsymbol{\varepsilon}$ and
inverse permeability $\boldsymbol{\mu}^{-1}$ are again uniaxial about
$\hat{\mathbf{n}}_*$. Here the outcome is the opposite of a new signal: the two
tensors obey the adjugate relation $\boldsymbol{\mu}^{-1} =
\operatorname{adj}(\boldsymbol{\varepsilon})$, which forces the two photon
polarizations to share a dispersion, so the gauge-sector vacuum \emph{birefringence
cancels}. What remains is a common-mode speed anisotropy---a direction-dependent
vacuum speed affecting both polarizations equally---whose observational status is
left open. That the two sectors inherit one axis from one geometry is the paper's
multi-channel concordance.

The gauge-sector story is developed analytically in the companion
Paper~VIII~\cite{menendez2026electricblock}, which derives the two blocks, proves the
adjugate cancellation theorem, and settles the electric sector that neither Paper~I
nor Paper~II supplied. The present paper is the independent, from-engine confirmation:
both blocks are read off the running lattice's hop holonomies, the cancellation is
verified numerically on those engine-extracted tensors, and the two papers are
intended as a back-to-back pair.

\subsection{Scope}
\label{subsec:scope}

In scope, and derived here: the closed-form kinematic speed split
(Sec.~\ref{sec:kinematic_channel}, \texttt{PASS}); the from-engine extraction of the
gauge-sector blocks and the null birefringence verdict
(Sec.~\ref{sec:gauge_sector}, \texttt{PASS}); and the structural axis-sharing of the
two channels. Explicitly deferred: the isotropic overall scale of the induced action
(the coupling $1/g^2$ Paper~I leaves open), so the gauge sector fixes anisotropic
\emph{structure} but not the absolute vacuum speed; the observational status of the
common-mode speed anisotropy (the single-domain-versus-$O_h$-restoration question,
\texttt{PART}); and the single-object \emph{dynamical} ($\operatorname{Tr}\ln T$)
extraction of the induced-action tensor, which is a stated open item
(Remark~\ref{rem:trlnt_open}). The quantitative falsifiability of the kinematic
prediction---its magnitude against current astrophysical bounds and its mapping onto
the Standard-Model Extension---is analyzed in the accompanying notes rather than
claimed in the prose. The audit table (Table~\ref{tab:audit}) is the canonical record
of which claims are \texttt{PASS}, \texttt{PART}, \texttt{STUB}, or \texttt{FAIL}.

% ============================================================
% audit_table.tex
%
% A "system audit" table mapping the paper's claims to the evidence
% backing each one.  This is *the* canonical source of truth for which
% claims are PASS / PART / STUB / FAIL --- prose elsewhere in the paper
% should not contradict it.
%
% Status semantics (paste into the abstract or front-matter as well):
%   PASS  -- derivation complete or experiment confirms the claim to
%            stated precision.
%   PART  -- mechanism demonstrated, full quantitative match pending;
%            specific gaps documented in the evidence column.
%   STUB  -- analytical result with no numerical confirmation yet,
%            or planned experiment not yet run.
%   FAIL  -- tested and disconfirmed.
%
% Use \texttt{} for code/experiment names and file paths inside cells.
% Long cells are fine -- the column widths are tuned for prose.
% ============================================================

\label{tab:audit}

{\scriptsize
\setlength{\tabcolsep}{4pt}
\renewcommand{\arraystretch}{0.95}
\begin{longtable}{@{}p{2.6cm}p{3.8cm}p{3.2cm}p{4.4cm}p{1.0cm}@{}}
\caption{\small System audit: paper claims mapped to supporting
derivations and experiments.  Status semantics in the front-matter
(\texttt{PASS} / \texttt{PART} / \texttt{STUB} / \texttt{FAIL}).}
\label{tab:audit_caption} \\
\hline
\textbf{Claim} & \textbf{Mechanism} & \textbf{Continuum / formal limit} & \textbf{Evidence} & \textbf{Status} \\
\hline
\endfirsthead

\multicolumn{5}{l}{\textit{(Table~\ref{tab:audit}, continued from previous page)}} \\
\hline
\textbf{Claim} & \textbf{Mechanism} & \textbf{Continuum / formal limit} & \textbf{Evidence} & \textbf{Status} \\
\hline
\endhead

\hline
\multicolumn{5}{r}{\textit{(continued on next page)}} \\
\endfoot

\hline
\endlastfoot

% ── Kinematic channel ───────────────────────────────────────────────
Kinematic birefringence along $(1,1,-1)$
    & Closed-form uniaxial split from the spinor propagator: $M_\text{eff} = \tfrac43(\mathbb{I}-\hat{\mathbf{n}}_*\hat{\mathbf{n}}_*^T)$, so $|H_\text{RGB}(\mathbf{k})|^2 = 1-\tfrac43|\mathbf{k}_\perp|^2$. Group-velocity split $\Delta v = \tfrac23\sin\omega\,|\mathbf{k}|$ (angular law $\sin\alpha$, massive sector) and $\omega$-independent dispersion-curvature split $4/3$ vs $0$ (photon + massive)
    & Isotropic $c$ under $O_h$ averaging / $a\to 0$
    & \emph{Analytic:} eigenphase dispersion $\tan\theta = \tan(\omega/2)/|H_\text{RGB}|$ and group velocity in closed form; $|H_\text{RGB}|^2 = 1$ exactly along $(1,1,-1)$ at all orders. \texttt{exp\_02\_closed\_form\_split\_check}: structure-factor and dispersion identities to $<10^{-12}$, speed coefficients to $<10^{-5}$. \emph{Observability:} \texttt{exp\_01\_optical\_axis\_isotropy} ($129^3$, both engines) -- isotropic pulse flattens perpendicular to $(1,1,-1)$ (ratio $\to 0.01$), flat-axis alignment $=1.0000$, $\omega$-independent, survives exact A=1 renormalization.
    & \texttt{PASS} \\

% ── Gauge sector ────────────────────────────────────────────────────
Gauge-sector birefringence ($\mathbf{Q}$ eigenvalues $\{4,4,16\}$): uniaxial structure derived; observability \emph{open}
    & Induced photon action $\tfrac{ca^4}{2}\mathbf{F}^T\mathbf{Q}\mathbf{F}$; special direction Hodge-duals to the same axis $(1,1,-1)$. The bipartite RGB/CMY structure breaks $O_h$ to a uniaxial subgroup about $(1,1,-1)$ ($\mathbf{Q}$ fixed by $12/48$, RGB by $6/48$), so the action is genuinely uniaxially anisotropic
    & $O_h$-averaged Maxwell is the \emph{orientation-averaged} form, not what a fixed lattice's observables see; quantitative $\mathbf{Q}$ role is the coupling ratio $g_3^2/g_2^2 = 3/2$ (Paper~II)
    & \emph{Structure derived} (Paper~I App.~B; Paper~II). \emph{Observability OPEN:} no symmetry forces isotropy (the $O_h$-average is an external imposition the uniaxial dynamics do not justify); by the \texttt{exp\_01} analogy the anisotropy may survive in observables, where read literally it is a dimension-4, $O(1)$ effect excluded by $\sim$30--37 orders. The one possible escape -- electric/magnetic cancellation in the full dispersion (Paper~I did the magnetic sector only) -- is uncomputed. \texttt{exp\_03} (full E+B Peierls dynamics) is decisive. See \texttt{notes/gauge\_sector\_structural\_conclusion.md}.
    & \texttt{PART} \\

% ── Concordance ─────────────────────────────────────────────────────
Multi-channel concordance (P9)
    & Kinematic and gauge sectors inherit the \emph{same} optical axis $(1,1,-1)$ from the common bipartite geometry ($M_\text{eff}$ and $\mathbf{Q}$ special directions coincide)
    & N/A (cross-sector structural identity)
    & Axis-sharing is a \emph{structural} identity (derived): $M_\text{eff}$ and $\mathbf{Q}$ both inherit $(1,1,-1)$ from the same bipartite geometry. The kinematic dispersion is the one frontier-testable observational channel; the gauge sector's observable status is itself \emph{open} (see its row -- it may be unobservable, suppressed, or an excluded $O(1)$ effect, pending \texttt{exp\_03}). P9's strength as a multi-witness empirical claim is therefore hostage to that outcome.
    & \texttt{PART} \\

\hline
\end{longtable}
}



\section{Kinematic Dispersion Anisotropy along $(1,1,-1)$}
% kinematic_channel.tex
% The kinematic directional-dispersion channel: closed-form axial/transverse
% group-speed contrast along the optical axis (1,1,-1).
% NOTE (terminology, external review 2026-07-17): "birefringence" is reserved
% for the gauge sector; this channel is directional dispersion anisotropy.
%
% STATUS: closed-form split DERIVED and numerically verified
% (notes/closed_form_speed_split.md; exp_02). One framing decision remains
% before flipping the audit row PART->PASS: which observable the
% "speed split" names (group-velocity vs dispersion-curvature) -- see the
% photon-limit remark below.  Cross-checks: exp_01 (observability), exp_02
% (analytic identities).

\label{sec:kinematic_channel}

A terminological caution frames this section. The effect below is a uniaxial
\emph{directional dispersion anisotropy}: it contrasts propagation along
$\hat{\mathbf{n}}_*$ with propagation across it. It is \emph{not} birefringence in
the optical sense---two polarization eigenmodes with different speeds at one fixed
wavevector---and we reserve that word for the gauge sector
(Sec.~\ref{sec:gauge_sector}). The two branches $\lambda_\pm(\mathbf{k})$ below are
the eigenvalue branches of the spinor transfer operator (the paired
positive/negative-frequency spinor modes), not two optical polarizations. A second
caution concerns interpretation: the transfer operator is \emph{non-unitary} away
from the axis ($|\lambda_\pm|^2<1$), so identifying the eigenphase gradient with a
physical group velocity needs an argument. The A=1 normalization supplies it for the
\emph{massive} sector but---as we show---\emph{not} in the massless limit, where the
eigenphase does not disperse at all (Sec.~\ref{subsec:a1_normalization},
Prop.~\ref{prop:massless_flat}). What is rigorous throughout is the exact
axial--transverse anisotropy of the structure-factor magnitude; the group-velocity
reading is earned for $\omega\neq0$, and the massless and photon readings are treated
with the care they require.

The kinematic channel inherits Paper~I's single-tick spinor
propagator~\cite{menendez2026geometryfirst}
\begin{equation}
    T(\mathbf{k}) = \begin{pmatrix}
        i\sin(\omega/2) & \cos(\omega/2)\,\tilde{H}_\text{RGB}(\mathbf{k}) \\
        \cos(\omega/2)\,\tilde{H}_\text{CMY}(\mathbf{k}) & i\sin(\omega/2)
    \end{pmatrix},
    \qquad
    \tilde{H}_\text{CMY} = \overline{\tilde{H}_\text{RGB}},
    \label{eq:propagator}
\end{equation}
whose eigenvalues
$\lambda_\pm(\mathbf{k}) = i\sin(\omega/2) \pm
\cos(\omega/2)\,|\tilde{H}_\text{RGB}(\mathbf{k})|$ carry the entire
$\mathbf{k}$-dependence through the single scalar
$|\tilde{H}_\text{RGB}|$.  Here $\omega$ is the per-tick Zitterbewegung
rate ($m = \sin(\omega/2)/a$), and the massless \emph{spinor} limit is
$\omega\to0$ (its identification with a photon is examined, and not asserted, in
Remark~\ref{rem:photon_limit}).  The RGB basis vectors are
$\mathbf{V}_1=(1,1,1)$, $\mathbf{V}_2=(1,-1,-1)$, $\mathbf{V}_3=(-1,1,-1)$,
with $\mathbf{V}_1+\mathbf{V}_2+\mathbf{V}_3=(1,1,-1)$ fixing the optical
axis $\hat{\mathbf{n}}_* = (1,1,-1)/\sqrt{3}$.

\subsection{Exact structure factor and the perpendicular projector}
\label{subsec:structure_factor}

With $\tilde{H}_\text{RGB}(\mathbf{k}) = \tfrac13\sum_{\mathbf{V}\in\text{RGB}}
e^{i\mathbf{k}\cdot\mathbf{V}}$ and the difference vectors
$\mathbf{V}_1-\mathbf{V}_2=(0,2,2)$, $\mathbf{V}_1-\mathbf{V}_3=(2,0,2)$,
$\mathbf{V}_2-\mathbf{V}_3=(2,-2,0)$, the modulus is closed-form to all
orders:
\begin{equation}
    |\tilde{H}_\text{RGB}(\mathbf{k})|^2
    = \tfrac19\Bigl[\,3
        + 2\cos 2(k_y{+}k_z)
        + 2\cos 2(k_x{+}k_z)
        + 2\cos 2(k_x{-}k_y)\,\Bigr].
    \label{eq:H2_exact}
\end{equation}
Along the optical axis $\mathbf{k}=t\,(1,1,-1)$ every argument vanishes, so
$|\tilde{H}_\text{RGB}|^2 = 1$ \emph{exactly for all $t$}: the dispersion
is rigorously flat along $\hat{\mathbf{n}}_*$.  Expanding to second order
recovers Paper~I's $M_\text{eff}$, which takes the transparent form
\begin{equation}
    |\tilde{H}_\text{RGB}(\mathbf{k})|^2 = 1 - \mathbf{k}^T M_\text{eff}\,\mathbf{k}
        + O(k^4),
    \qquad
    M_\text{eff} = \tfrac43\bigl(\mathbb{I} - \hat{\mathbf{n}}_*\hat{\mathbf{n}}_*^{T}\bigr)
        = \tfrac43 P_\perp .
    \label{eq:Meff_projector}
\end{equation}
That is, $M_\text{eff}$ is $\tfrac43$ times the projector onto the plane
perpendicular to the optical axis (eigenvalues $\{4/3,4/3,0\}$).  The
dispersion therefore depends \emph{only} on the perpendicular wavevector
$\mathbf{k}_\perp = P_\perp\mathbf{k}$:
\begin{equation}
    |\tilde{H}_\text{RGB}(\mathbf{k})|^2 = 1 - \tfrac43|\mathbf{k}_\perp|^2 + O(k^4).
    \label{eq:H2_perp}
\end{equation}

\subsection{The eigenphase dispersion and the speed split}
\label{subsec:speed_split}

The mode energy (phase advance per tick) is
$\theta(\mathbf{k}) = \arg\lambda_+(\mathbf{k})$.  From
Eq.~\eqref{eq:propagator}, $\mathrm{Re}\,\lambda_+
= \cos(\omega/2)|\tilde{H}_\text{RGB}|$ and
$\mathrm{Im}\,\lambda_+ = \sin(\omega/2)$, giving the exact dispersion
\begin{equation}
    \tan\theta(\mathbf{k}) = \frac{\tan(\omega/2)}{|\tilde{H}_\text{RGB}(\mathbf{k})|}.
    \label{eq:dispersion_exact}
\end{equation}
At $\mathbf{k}=0$ this returns $\theta=\omega/2$ (the rest-mass phase);
off-axis $|\tilde{H}_\text{RGB}|<1$ raises $\theta$.  Differentiating
gives the group velocity in closed form,
\begin{equation}
    \mathbf{v}_g(\mathbf{k}) = \nabla_{\mathbf{k}}\theta
    = -\,\frac{\tfrac12\sin\omega\;\nabla|\tilde{H}_\text{RGB}|}
             {\cos^2(\omega/2)\,|\tilde{H}_\text{RGB}|^2 + \sin^2(\omega/2)},
    \label{eq:group_velocity}
\end{equation}
which, using Eq.~\eqref{eq:H2_perp}, reduces at small $\mathbf{k}$ to
\begin{equation}
    \mathbf{v}_g(\mathbf{k}) = \tfrac23\sin\omega\;\mathbf{k}_\perp + O(k^3).
    \label{eq:vg_small_k}
\end{equation}

\begin{proposition}[Axial--transverse group-velocity contrast]
\label{prop:speed_split}
For a wavevector at angle $\alpha$ to the optical axis
($|\mathbf{k}_\perp| = |\mathbf{k}|\sin\alpha$), the kinematic group speed
(read from the eigenphase, subject to the interpretation caveat below) is, to
leading order,
\begin{equation}
    |\mathbf{v}_g| = \tfrac23\sin\omega\;|\mathbf{k}|\,\sin\alpha .
    \label{eq:speed_law}
\end{equation}
The transverse direction ($\mathbf{k}\perp\hat{\mathbf{n}}_*$, $\alpha=90^\circ$)
carries the maximal group speed $\tfrac23\sin\omega\,|\mathbf{k}|$; the axial
direction ($\mathbf{k}\parallel\hat{\mathbf{n}}_*$, $\alpha=0$) is frozen,
$|\mathbf{v}_g|=0$.  The axial--transverse contrast is
\begin{equation}
    \Delta v \equiv v_\perp - v_\parallel = \tfrac23\sin\omega\;|\mathbf{k}|,
    \label{eq:split}
\end{equation}
with $\hat{\mathbf{n}}_* = (1,1,-1)/\sqrt{3}$ the frozen (axial) direction.
This is a contrast between two propagation \emph{directions}, not a
polarization split at fixed $\mathbf{k}$.
\end{proposition}

\subsection{The A=1 normalization and the massless limit}
\label{subsec:a1_normalization}

Away from the axis the transfer operator is not unitary:
$|\lambda_\pm(\mathbf{k})|^2 = \sin^2(\omega/2) + \cos^2(\omega/2)\,
|\tilde H_\text{RGB}(\mathbf{k})|^2 < 1$, so each eigenvalue carries amplitude
change as well as phase, and $\nabla_{\mathbf{k}}\theta$ is not \emph{a priori} a
physical group velocity. The A=1 rule bears on this, with a caveat we state carefully.
The scheduler renormalizes the amplitude to unit norm every tick---a real, positive
rescaling $\psi\to\psi/\|\psi\|$ that leaves every phase untouched and writes back
$\phi=\arg\psi$. For a pure Bloch \emph{eigenmode} this removes the scalar attenuation
while preserving the eigenphase, replacing $\lambda_\pm$ by
$\lambda_\pm/|\lambda_\pm|$, so $\theta_\pm(\mathbf{k}) = \arg\lambda_\pm(\mathbf{k})$
is the physical single-mode dispersion. For a general state the renormalization is
\emph{nonlinear}---a single global norm does not equalize the different $|\lambda(\mathbf{k})|$
of a packet's several $\mathbf{k}$-components---so it is \emph{not} a unitary operator on
the Hilbert space, and the group-velocity reading must be justified at the level of the
wave packet rather than by calling the map unitary. For a narrow-band packet a
stationary-phase estimate gives the centroid velocity as the phase-gradient
$\nabla_{\mathbf{k}}\theta$ (Eqs.~\eqref{eq:dispersion_exact},
\eqref{eq:group_velocity}), while the momentum dependence of $|\lambda|$ contributes an
additional amplitude filtering/spreading term, treated separately below and at $\omega=0$.
This transport reading is corroborated by \texttt{exp\_01} under the full nonlinear
dynamics; the distinction between the massive-sector phase dispersion and the massless
magnitude filtering is exactly the one drawn next.

\begin{proposition}[Massless-limit degeneracy]
\label{prop:massless_flat}
At $\omega=0$ the eigenvalues $\lambda_\pm = \pm|\tilde H_\text{RGB}(\mathbf{k})|$ are
real, so $\theta_\pm(\mathbf{k}) \in \{0,\pi\}$ independently of $\mathbf{k}$: the
eigenphase is exactly flat and the group velocity vanishes in every direction. (This
is immediate from the closed form---$\lambda_\pm$ real makes $\arg\lambda_\pm$
constant---and a direct finite-difference sweep confirms
$\max_{\mathbf{k}}|\nabla_{\mathbf{k}}\theta_+| = 0$ at $\omega=0$, versus $O(1)$ for
$\omega\neq0$.)
\end{proposition}

The A=1 normalization does not restore a dispersion here---the phase is already
$\mathbf{k}$-independent---so at $\omega=0$ the entire anisotropy is carried by the
eigenvalue \emph{magnitude} $|\tilde H_\text{RGB}|$, and there is no eigenphase energy
dispersion. It is important not to over-read the single-mode step: fixing
$|\lambda_{\mathbf{k}}|\to1$ per plane wave would flatten one mode's magnitude, but the
physical A=1 rule renormalizes the \emph{total} norm of a state, which does \emph{not}
equalize the relative magnitudes across the $\mathbf{k}$-components of a wavepacket. The
magnitude anisotropy therefore survives into observable amplitudes---this is exactly the
$\omega$-independent anisotropic spreading that \texttt{exp\_01} measures (present and
identical at $\omega=0$). What is absent at $\omega=0$ is an energy \emph{dispersion},
not the anisotropy itself. Two consequences follow, and they set the honest scope of
the channel. First, the directional-dispersion (group-velocity) reading of this section
is a \emph{massive-sector} ($\omega\neq0$) effect, whereas the $\omega$-independent
magnitude anisotropy is a filtering/spreading effect, not an energy dispersion. Second,
a \emph{massless} energy dispersion---the object an astrophysical photon test would
constrain---is \emph{not} delivered by the $\omega\to0$ limit of the spinor propagator;
on this lattice the photon is the induced gauge field of Sec.~\ref{sec:gauge_sector},
whose
dispersion is the common-mode $\omega^2 = \mathbf{k}^{\!\top}\boldsymbol{\varepsilon}
\,\mathbf{k}$. This is the same reason the massless spinor mode is not, without a
further argument, a photon (Remark~\ref{rem:photon_limit}).

\begin{remark}[The nonlinear renormalization carries a foundational constraint]
\label{rem:nonlinear_qm}
Because the A=1 rescaling is a norm-restoring \emph{nonlinear} evolution, it inherits a
constraint that is independent of Lorentz tests. Norm-restoring nonlinear modifications of
quantum mechanics are of the postselection type and generically permit superluminal
signaling between entangled parties~\cite{weinberg1989,gisin1990,polchinski1991}:
concretely, for an entangled pair whose $A$-side branches differ in $|\mathbf{k}_\perp|$,
the universal differential damping of Sec.~\ref{subsec:as_constructed} together with a
\emph{global} renormalization can change $B$'s marginals depending on $A$'s local basis
choice, which a local completely-positive trace-preserving (CPTP) map cannot do. There is
also an internal tension to square: were the renormalization a harmless gauge, the
$\omega=0$ filtering would be unobservable, yet \texttt{exp\_01} and the analysis above
insist it \emph{is} observable. We do not resolve this here; we flag the three routes that
would, each with different phenomenology. \emph{(a)} Implement A=1 as a genuine CPTP map by
dilation, so the lost norm flows to an explicit sink sector---anisotropic \emph{absorption}
into residue rather than a nonlinear rescale, which signals nothing but opens a decay
channel with its own bounds (this dovetails the program's ``nothingness'' ontology).
\emph{(b)} A minimum-probability floor $\delta p_\text{min}$ that cuts off the differential
damping (the same regulator invoked for the filtering). \emph{(c)} An argument that the
A=1 rule, restricted to the physically preparable state space, never generates the
dangerous entangled configurations. This touches every channel that leans on the
renormalization and warrants a dedicated treatment; it is a stated open foundational item.
\end{remark}

\subsection{The as-constructed reading: an $O(1)$ matter-sector exclusion}
\label{subsec:as_constructed}

The closed forms above are exact, and it is tempting to read the resulting
anisotropy as only the small, dimension-six fractional energy shift of
Sec.~\ref{sec:falsifiability}. Taken literally, though, the bare single-tick
spinor forces three much stronger, order-unity statements about the matter sector,
all with the same geometric origin: the difference vectors
$\mathbf{V}_i-\mathbf{V}_j$ of Sec.~\ref{subsec:structure_factor} all lie in the
plane perpendicular to $\hat{\mathbf{n}}_*$, so \emph{every} axial dependence of
$\tilde H_\text{RGB}$ is a pure global phase that the modulus discards. We state
them plainly; all are verified in \texttt{exp\_05\_matter\_sector\_order}.

\begin{proposition}[Exact axial flat band]
\label{prop:flat_band}
$\theta(\mathbf{k}) = \theta(\mathbf{k}_\perp)$ to all orders, so
$\partial\theta/\partial k_\parallel \equiv 0$ at \emph{every} $\mathbf{k}$, not
merely at small $\mathbf{k}$. A wavepacket with momentum along $\hat{\mathbf{n}}_*$
does not propagate, and---since $\dot{x}_\parallel = \partial\theta/\partial
k_\parallel$ and the two-component Berry connection contributes no anomalous
velocity---$x_\parallel$ stays inert even under forces that change $k_\parallel$.
This is an exact flat band along the axis, with the attendant infinite degeneracy;
a hydrogen-like state built on this kinetic term is dynamically two-dimensional.
\end{proposition}

Two further consequences follow at fixed $|\mathbf{k}|$. First, the kinetic energy
is \emph{entirely} directional:
\begin{equation}
    \theta(\mathbf{k}) - \tfrac{\omega}{2}
      = \tfrac13\sin\omega\,|\mathbf{k}|^2\sin^2\!\alpha + O(k^4),
    \label{eq:kinetic_energy}
\end{equation}
modulated by $\sin^2\alpha$ with \emph{no} isotropic piece for the anisotropy to be
a small correction to. The dimension-six fractional number of
Sec.~\ref{sec:falsifiability} is the energy anisotropy relative to the \emph{rest}
phase $\omega/2$; the quantity a Hughes--Drever--class experiment compares---kinetic
and level structure versus orientation---is anisotropic at order unity. If the
spinor is any laboratory particle, this is excluded at zeroth order, no precision
required. Second, the scale cannot be tuned away: the lattice perpendicular kinetic
coefficient over the continuum $\mathbf{k}^2/2m$ (with $m=\sin(\omega/2)/a$) is
$\tfrac43\sin^2(\omega/2)\cos(\omega/2)\le 8/(9\sqrt3)\approx0.513$ at \emph{any}
$\omega$, so the single-tick spinor is simultaneously axially inert, $O(1)$
anisotropic, and kinetically $(ma)^2$-suppressed relative to a continuum particle.

\begin{proposition}[Linearly accumulating $\omega=0$ filtering]
\label{prop:filtering}
The eigenvalue magnitude obeys $|\lambda_+(\mathbf{k})|^2 = \sin^2(\omega/2) +
\cos^2(\omega/2)\,|\tilde H_\text{RGB}(\mathbf{k})|^2$---unity along the axis, and
$1-\tfrac43\cos^2(\omega/2)|\mathbf{k}_\perp|^2$ across it. The per-tick intensity
damping of a transverse relative to an axial component is therefore
$\tfrac43\cos^2(\omega/2)(k_\perp a)^2$; multiplied by the tick rate $c/a$, the
physical rate is $\tfrac43\cos^2(\omega/2)\,k_\perp^2\,a\,c$---\emph{linear} in $a$,
because the dimension-six per-tick suppression is repaid by the $a^{-1}$ tick count.
The dispersion channels do not accumulate (a velocity shift stays $(ka)^2$ forever);
this amplitude channel does.
\end{proposition}

At $a=\ell_P$ the resulting $e$-folding times are days to minutes---an electron at
hydrogenic momentum $\approx5$ days, a thermal $\mathrm{N}_2$ molecule
$\approx49$ minutes, a cold neutron $\approx7$ days. Even allowing for collisional
rethermalization (a gas relaxes to a steady-state momentum-distribution anisotropy
of order the filter-to-collision-rate ratio) and generous modeling slack, the effect
starts tens of orders of magnitude above co-magnetometer and matter-interferometer
bounds. The ``testability open'' status this channel formerly carried is therefore
too sanguine.

\paragraph{The honest verdict, and the escape.} Read as literal laboratory matter,
the bare single-tick spinor is \emph{excluded as constructed}: axially inert
(Prop.~\ref{prop:flat_band}), $O(1)$ anisotropic in its kinetic sector, kinematically
$(ma)^2$-mismatched to a continuum particle, and amplitude-filtering
(Prop.~\ref{prop:filtering}) at rates far above bound. This is the matter-sector
counterpart of the gauge-sector common-mode tension (Sec.~\ref{sec:gauge_sector}),
and it has the same three-of-four geometric origin. The escape is the one the program
already owns: physical matter is \emph{emergent}---composite ``session'' excitations
with their own dispersion, not the bare single-tick spinor---so the closed forms above
are the \emph{substrate} kinematics, and the observable matter dispersion (with it the
sidereal channel of Sec.~\ref{sec:falsifiability}) follows only once the
token$\to$matter map is exhibited; the program's minimum-probability floor
$\delta p_\text{min}$ is the natural regulator that halts the filtering. Until that map
is delivered, the matter-sector channels read literally are \texttt{FAIL} as
constructed, with the emergence map the open route to viability
(Table~\ref{tab:audit}). What changes is \emph{not} the closed-form
split~\eqref{eq:split}, which remains exact as substrate kinematics, but only the
claim that it registers as a small observable effect on a laboratory particle.

Equations~\eqref{eq:H2_exact}, \eqref{eq:dispersion_exact},
\eqref{eq:group_velocity} and Proposition~\ref{prop:speed_split} have been
checked against independent direct computation in
\texttt{exp\_02\_closed\_form\_split\_check}: the structure-factor and
dispersion identities to floating-point precision ($<10^{-12}$, versus the
triple sum and versus $\arg\lambda_+$ respectively), and the group
velocity and principal speed coefficients $\tfrac23\sin\omega$ and $0$ to
finite-difference truncation ($<10^{-5}$); the order-unity consequences of
Sec.~\ref{subsec:as_constructed} are verified in
\texttt{exp\_05\_matter\_sector\_order}. These checks establish the closed-form
identities (the kinematic row's \emph{Analytic} content is \texttt{PASS}); the
\emph{observable} reading of the same identities on a laboratory particle is the
as-constructed \texttt{FAIL} of Sec.~\ref{subsec:as_constructed}, pending the
emergence map. The
observability of the underlying anisotropy in renormalized amplitudes was
established separately in \texttt{exp\_01} (Table~\ref{tab:audit};
\texttt{notes/optical\_axis\_renormalization\_test.md}).

\begin{remark}[The massless mode is not automatically a photon]
\label{rem:photon_limit}
Proposition~\ref{prop:massless_flat} already shows that the $\omega\to0$ eigenphase
does not disperse. There is a second, independent reason not to read the massless
limit as a photon result: the $\omega=0$ limit of a spinor propagator is, in the
first instance, a massless \emph{spinor} mode, and identifying the two-component mode
with a photon---its helicity content, gauge character, and representation---is a
separate step this paper does not take. The robust framework-internal content of this
section is thus the closed-form uniaxial anisotropy of $|\tilde H_\text{RGB}|$ about
$\hat{\mathbf{n}}_*$ and its survival under A=1 renormalization (\texttt{exp\_01}),
with the massive-sector group velocity~\eqref{eq:split} as its clearest observable
instance; the photon-sector dispersion is treated separately in
Sec.~\ref{sec:gauge_sector}.
\end{remark}

% ── Scoping remark: mode count is unchanged by the split ──────────────
\begin{remark}[The split adds no propagating modes]
\label{rem:no_new_modes}
The axial and transverse group speeds derived above are the two
\emph{direction-dependent phases} of the bipartite propagator's existing
eigenvalue branches $\lambda_\pm(\mathbf{k})
= i\sin(\omega/2) \pm \cos(\omega/2)\,|\tilde{H}_\text{RGB}(\mathbf{k})|$.
The entire directional anisotropy enters through the single scalar
$|\tilde{H}_\text{RGB}(\mathbf{k})|$, whose low-$\mathbf{k}$ expansion
$1 - \tfrac43|\mathbf{k}_\perp|^2$ is flat along $\hat{\mathbf{n}}_*$ and
curved (coefficient $4/3$) in the perpendicular plane.  The closed-form
split is thus a modulation of the \emph{same} two branches---the
$(\psi_R,\psi_L)$ pair of the continuum Dirac spinor---and introduces no
additional branch: the mode count is two, before and after.

This is logically independent of the lattice's fermion-doubling content.
Doublers are a Brillouin-zone-corner phenomenon, whereas the
axial--transverse split is the second-order term of an expansion
about the physical mode at $\mathbf{k}=0$; the two occupy disjoint regions
of $\mathbf{k}$-space.  Paper~I deliberately defers the doubler accounting
for the bipartite $T_\text{diamond}$ Dirac operator (Paper~I~\S6.4,
\cite{menendez2026geometryfirst}) to a separate one-loop
lattice-perturbation-theory calculation,
and nothing in the present derivation depends on, alters, or settles that
count.  This $\mathbf{k}$-space disjointness argument is inherited from
Paper~I's framework and is not a new result of the present paper; in
particular, the optical-axis dispersion anisotropy does \emph{not} imply a larger
particle zoo.
\end{remark}

% ── Scoping remark: directional dispersion and packet spreading are the ──
% ── same channel read two ways, not two new falsifiers.  Folds the ────
% ── proposed post-v1 "Experiment 5" into this section per the program ─
% ── decision (4 falsifiers + 1 demonstrator; Exp 5 is the demonstrator).
\begin{remark}[Directional dispersion and wave-packet spreading restate this channel]
\label{rem:spreading_same_channel}
Two real-space signatures sometimes proposed as independent tests are, on
this lattice, the \emph{same} kinematic anisotropy read through different
observables and add no new claim. First, a direction-dependent propagation
speed---sometimes labeled ``vacuum birefringence'' in the literature---is, on
this lattice, precisely the directional dispersion anisotropy of
Proposition~\ref{prop:speed_split} and Eq.~\eqref{eq:H2_perp}, not a distinct
fixed-$\mathbf{k}$ polarization split; it is this section's subject, not a
separate prediction. Second, the anisotropic \emph{spreading} of an initially
isotropic packet is the real-space dual of the $\omega$-independent
structure-factor magnitude curvature $|\tilde H_\text{RGB}|^2 = 1 -
\tfrac43|\mathbf{k}_\perp|^2$ (Eq.~\eqref{eq:H2_perp})---\emph{not} the
massive-sector group velocity of Eq.~\eqref{eq:group_velocity}, which vanishes at
$\omega=0$---and is exactly the observable that \texttt{exp\_01} already measures
(the pulse flattens perpendicular to $\hat{\mathbf{n}}_*$, present and identical at
$\omega=0$). That this spreading is $\omega$-independent, whereas the group-velocity
contrast~\eqref{eq:split} scales as $\sin\omega$, is precisely the distinction of
Sec.~\ref{subsec:a1_normalization}: the magnitude anisotropy is a filtering/spreading
effect, the eigenphase group velocity a massive-sector energy dispersion. The isotropic-spread rate itself
($\mathrm{d}^2E/\mathrm{d}p^2$, direction-averaged) is the province of the
Lorentz-violating dispersion program, whose real-space demonstrator lives
there rather than here; only its \emph{$(1,1,-1)$-aligned anisotropy} is
this paper's result. Any claim that the lattice packet spreads ``faster than
continuum QFT'' must be stated as the residual that survives the continuum
limit $a\to0$: a free packet already spreads, and a lattice excess that
vanishes as $a\to0$ is numerical dispersion, not physics. What survives
$a\to0$ here is the \emph{direction dependence} fixed by
$\hat{\mathbf{n}}_*$, not an isotropic spread-rate enhancement.
\end{remark}


\section{Gauge-Sector Response and the Birefringence Cancellation along $(1,1,-1)$}
% gauge_sector.tex
% The gauge-sector channel: the induced Maxwell action's two anisotropic blocks
% (permittivity eps=P={1,4,4}, inverse permeability mu^-1=Q_B={4,4,16}), the
% adjugate relation, the null birefringence verdict, and the (separate) common-mode
% vacuum-speed anisotropy.
%
% STATUS (audit_table.tex is authority):
%   birefringence verdict (cancels) -> PASS  (null split; engine blocks + solver)
%   common-mode vacuum-speed anisotropy -> PART  (single-domain factor ~2 vs O_h)
% Framing: geometric STRUCTURE (from-engine holonomy) x scalar MAGNITUDE (deferred
% 1/g^2). The single-object dynamical (Tr ln T) tensor extraction is a STATED OPEN
% item (see the remark at the end). Derivation + adjugate theorem: Paper~VIII
% \cite{menendez2026electricblock}. Engine extraction + verdict: exp_03_geometric_blocks,
% exp_03_dispersion_core; notes/exp_03_dynamical_confirmation_plan.md.

\label{sec:gauge_sector}

The same bipartite geometry that fixes the kinematic optical axis also fixes a
\emph{second} anisotropy, in the gauge sector.  Integrating the matter tokens out
of the A=1 dynamics in a background electromagnetic field induces an effective
Maxwell action whose energy density is the anisotropic quadratic form
\begin{equation}
    u = \tfrac12\bigl(\mathbf{E}^{T}\boldsymbol{\varepsilon}\,\mathbf{E}
        + \mathbf{B}^{T}\boldsymbol{\mu}^{-1}\mathbf{B}\bigr),
    \label{eq:induced_action}
\end{equation}
with a permittivity tensor $\boldsymbol{\varepsilon}$ and an inverse-permeability
tensor $\boldsymbol{\mu}^{-1}$ that are \emph{not} isotropic.  Both are uniaxial
about the same optical axis $\hat{\mathbf{n}}_* = (1,1,-1)/\sqrt3$ that governs the
kinematic channel (Sec.~\ref{sec:kinematic_channel}).  We report the two tensors as
they are read directly off the engine's hop geometry, establish that the induced
birefringence \emph{cancels}, and separate that (settled) null result from the
distinct---and open---question of the common-mode vacuum speed.  The overall scale
of Eq.~\eqref{eq:induced_action} carries the one undetermined coupling $1/g^2$ that
Paper~I leaves open; throughout, we display anisotropic \emph{structure} and treat
that isotropic \emph{magnitude} as deferred.  The analytic derivation of the blocks
and of the cancellation theorem is the companion Paper~VIII
\cite{menendez2026electricblock}; the present section is the independent from-engine
confirmation.

\subsection{The two blocks from the hop geometry}
\label{subsec:blocks}

Write the three RGB hop vectors as $\mathbf{V}_1=(1,1,1)$, $\mathbf{V}_2=(1,-1,-1)$,
$\mathbf{V}_3=(-1,1,-1)$ (as in Sec.~\ref{sec:kinematic_channel}).  The two blocks
are the two natural quadratic forms these vectors build:
\begin{equation}
    \boldsymbol{\varepsilon} = P = \sum_{a} \mathbf{V}_a\mathbf{V}_a^{T},
    \qquad
    \boldsymbol{\mu}^{-1} = Q_B
        = \sum_{a<b} (\mathbf{V}_a\times\mathbf{V}_b)(\mathbf{V}_a\times\mathbf{V}_b)^{T}.
    \label{eq:blocks_def}
\end{equation}
The electric block sums the hop vectors themselves; the magnetic block sums their
oriented \emph{plaquette} areas $\mathbf{V}_a\times\mathbf{V}_b$.  Explicitly,
\begin{equation}
    \boldsymbol{\varepsilon} = \begin{pmatrix} 3&-1&1\\-1&3&1\\1&1&3 \end{pmatrix},
    \quad \text{eig}\{1,4,4\};
    \qquad
    \boldsymbol{\mu}^{-1} = \begin{pmatrix} 8&4&-4\\4&8&-4\\-4&-4&8 \end{pmatrix},
    \quad \text{eig}\{4,4,16\}.
    \label{eq:blocks_explicit}
\end{equation}
Both share the eigenvector $\hat{\mathbf{n}}_*=(1,1,-1)/\sqrt3$, but with opposite
sense: the optical axis is the \emph{suppressed} ($\varepsilon=1$) direction of the
permittivity and the \emph{enhanced} ($\mu^{-1}=16$) direction of the inverse
permeability.  The origin of the axis is geometric: $\hat{\mathbf{n}}_*$ is the one
cube body-diagonal \emph{absent} from the hop set $\{\mathbf{V}_a\}$, so the hop set
breaks the cubic group $O_h$ to the trigonal $D_{3d}$ about $\hat{\mathbf{n}}_*$
\cite{menendez2026electricblock}.

Each block is read off the running engine, not asserted.  The magnetic block is the
spatial plaquette holonomy of a uniform field: with the Peierls link phase
$\exp(i\,\tfrac12(\mathbf{A}(\mathbf{x})+\mathbf{A}(\mathbf{x}-\mathbf{V}))\cdot\mathbf{V})$
threaded through the hop, the smallest closed loop spanned by $\mathbf{V}_a,\mathbf{V}_b$
carries flux $(\mathbf{V}_a\times\mathbf{V}_b)\cdot\mathbf{B}$, and summing its square
over the RGB planes reproduces $Q_B$ (the dcl-core \texttt{exp\_04} route).  The
electric block has no purely spatial holonomy---the field enters the tick rule
\emph{on-site}, through $\delta\phi(\mathbf{x}) = \omega + V(\mathbf{x})$ with
$V=-\mathbf{E}\cdot\mathbf{x}$---so it is read from the \emph{temporal} plaquette:
seeding a uniform real state and taking one real step recovers
$\delta\phi = \omega + V$ from $\mathrm{Im}\,\psi_R^{\text{new}}/\psi_R
= \sin(\delta\phi/2)$, and the loop
$\Theta_a(\mathbf{x}) = \delta\phi(\mathbf{x}) - \delta\phi(\mathbf{x}+\mathbf{V}_a)
= \mathbf{V}_a\cdot\mathbf{E}$ (the mass term $\omega$ cancels around the loop)
gives $\sum_a\Theta_a^2 = \mathbf{E}^{T}P\,\mathbf{E}$
\cite{menendez2026electricblock}.  The experiment
\texttt{exp\_03\_geometric\_blocks} performs both extractions from one engine and
recovers Eq.~\eqref{eq:blocks_explicit} with the optical axis in the expected
(suppressed/enhanced) sense; the on-site coupling reproduces $\omega+V$ to
$4\times10^{-19}$ (unit tick weight), and the extracted $P$ is independent of
$\omega$ to $7\times10^{-15}$ (the mass cancels), so the number genuinely exercises
\texttt{HopOperator.step}.

\subsection{The covariant relation and the null birefringence verdict}
\label{subsec:verdict}

The two blocks are not independent.  For \emph{any} three vectors, the sum of
squared plaquette areas is the adjugate of the sum of squared vectors,
\begin{equation}
    Q_B = \sum_{a<b}(\mathbf{V}_a\times\mathbf{V}_b)(\mathbf{V}_a\times\mathbf{V}_b)^{T}
        = \operatorname{adj}\!\Bigl(\sum_a\mathbf{V}_a\mathbf{V}_a^{T}\Bigr)
        = \operatorname{adj}(P),
    \label{eq:adjugate}
\end{equation}
so $\boldsymbol{\mu}^{-1} = \operatorname{adj}(\boldsymbol{\varepsilon})$ is a
structural identity, not a coincidence of the octahedral lattice
\cite{menendez2026electricblock}.  The two engine-extracted blocks satisfy it: they
commute to $4\times10^{-15}$ and $\operatorname{adj}(\boldsymbol{\varepsilon})$
matches $\boldsymbol{\mu}^{-1}$ to $5\times10^{-15}$
(\texttt{exp\_03\_geometric\_blocks}).  Equation~\eqref{eq:adjugate} is exactly the
condition that removes the birefringence.

\begin{proposition}[Null gauge-sector birefringence]
\label{prop:null_split}
For a medium with $\boldsymbol{\mu}^{-1} = \operatorname{adj}(\boldsymbol{\varepsilon})$,
the photon dispersion factorizes as
\begin{equation}
    \det\!\bigl(K\,\boldsymbol{\mu}^{-1}K + \omega^2\boldsymbol{\varepsilon}\bigr)
    = \det(\boldsymbol{\varepsilon})\,\omega^2\,
      \bigl(\omega^2 - \mathbf{k}^{T}\boldsymbol{\varepsilon}\,\mathbf{k}\bigr)^2 ,
    \qquad K = [\hat{\mathbf{k}}]_\times ,
    \label{eq:dispersion_factored}
\end{equation}
a \emph{double} transverse root for every propagation direction.  Both
polarizations share $\omega^2 = \mathbf{k}^{T}\boldsymbol{\varepsilon}\,\mathbf{k}$,
so the split
$\Delta(\omega^2)\propto |\varepsilon_a\mu^{-1}_a - \varepsilon_p\mu^{-1}_p|
= |\det\boldsymbol{\varepsilon}-\det\boldsymbol{\varepsilon}| = 0$
vanishes.  The cancellation is invariant under independent rescaling of
$\boldsymbol{\varepsilon}$ and $\boldsymbol{\mu}^{-1}$, hence immune to the deferred
$1/g^2$, and holds for any hop vectors \textup{\cite{menendez2026electricblock}}.
\end{proposition}

Feeding the two \emph{engine-extracted} blocks of Eq.~\eqref{eq:blocks_explicit}
into the plane-wave dispersion of Proposition~\ref{prop:null_split} confirms this
numerically: the experiment \texttt{exp\_03\_dispersion\_core} solves the
generalized eigenproblem $-K\boldsymbol{\mu}^{-1}K\,\mathbf{E}
= v^2\boldsymbol{\varepsilon}\,\mathbf{E}$ and finds the two transverse speeds equal
to $\max|v_+-v_-| = 1.6\times10^{-15}$ over $2000$ directions, while a deliberately
detuned control (breaking Eq.~\eqref{eq:adjugate}) returns a split of $0.5$---so the
null is a genuine cancellation, not an insensitive estimator.  This discharges the
gauge-sector \emph{birefringence} audit row to \texttt{PASS}
(Table~\ref{tab:audit}).  Independently, the same operator anisotropy \emph{survives}
the A=1 token renormalization: driving the engine's full multi-tick dynamics with a
background field, the induced second-order response is uniaxial about
$\hat{\mathbf{n}}_*$ with the correct opposite senses ($\mathbf{B}$ axis-enhanced,
$\mathbf{E}$ axis-suppressed), so the cancellation is a property of the renormalized
observables, not only of the tree-level operator.

\subsection{The common-mode vacuum speed}
\label{subsec:common_mode}

The null split does \emph{not} make the gauge sector isotropic.  The shared
dispersion of Proposition~\ref{prop:null_split},
$\omega^2 = \mathbf{k}^{T}\boldsymbol{\varepsilon}\,\mathbf{k}$, is
direction-dependent: the common-mode phase speed
$v^2(\hat{\mathbf{k}}) = \hat{\mathbf{k}}^{T}\boldsymbol{\varepsilon}\,\hat{\mathbf{k}}$
ranges over the eigenvalues of $\boldsymbol{\varepsilon}$, i.e.\ $v^2\in[1,4]$---a
factor-$\sim2$ anisotropy of the vacuum speed about $\hat{\mathbf{n}}_*$, affecting
\emph{both} polarizations equally.  The dispersion solver on the engine blocks
returns $v\in[1.00,2.00]$ for a single lattice.  This is a speed anisotropy, not
birefringence, and its observational status is genuinely open.

\begin{remark}[Single domain versus $O_h$ restoration]
\label{rem:oh_restoration}
The factor-$\sim2$ anisotropy is the response of a \emph{single} trigonal domain.
Averaging the blocks over the four ways of choosing three of the four cube
body-diagonals as hop directions gives
$\langle\boldsymbol{\varepsilon}\rangle = 3\mathbb{I}$ and
$\langle\boldsymbol{\mu}^{-1}\rangle = 8\mathbb{I}$, restoring $O_h$ and removing the
anisotropy entirely (\texttt{exp\_03\_geometric\_blocks}: residual $10^{-16}$).  But
the physical lattice has \emph{one fixed} hop set---its optical axis is a specific
body-diagonal---so this four-domain average is a fictitious ensemble the lattice does
not realize, and there is no dynamical mechanism that performs it.  Whether the
common-mode anisotropy is therefore a real prediction (single domain) or is washed
out (restored) is the isotropy-restoration question this paper leaves at
\texttt{PART}; it is shared with the companion Paper~VIII
\cite{menendez2026electricblock} and is distinct from the null birefringence, which
holds in \emph{every} domain via Eq.~\eqref{eq:adjugate} regardless of the answer.
\end{remark}

\subsection{Scope: structure, magnitude, and the deferred dynamical tensor}
\label{subsec:gauge_scope}

The results above are the anisotropic \emph{structure} of the induced action
(Eq.~\eqref{eq:blocks_explicit}), computed from the engine hop geometry, together
with the exact cancellation it implies.  The isotropic \emph{magnitude}---the single
scalar $1/g^2$ multiplying Eq.~\eqref{eq:induced_action}---is the coupling Paper~I
leaves undetermined and is not a prediction of the vacuum speed here.  Because the
birefringence verdict is prefactor-independent (Proposition~\ref{prop:null_split}),
that omission does not weaken it.

\begin{remark}[The dynamical tensor extraction is a stated open item]
\label{rem:trlnt_open}
The blocks of Eq.~\eqref{eq:blocks_def} are obtained from the hop \emph{geometry}
(the plaquette holonomies), which is the tree-level induced action and is exact at
any lattice size.  One can also ask for the tensors as a genuinely \emph{dynamical},
momentum-space fermion-loop sum---the second-order response of the physical band of
the one-period transfer operator $T(\mathbf{k})$ summed over the Brillouin zone
($\operatorname{Tr}\ln T$).  We do not carry that computation out here, and we do not
claim it.  A naive band sum is not merely hard but Ward-\emph{divergent}: along the
optical axis the physical band is exactly flat (all three hops project equally onto
$\hat{\mathbf{n}}_*$, the $D_{3d}$ symmetry), so the intraband denominator
$\lambda(\mathbf{k})-\lambda(\mathbf{k}\pm\mathbf{q})$ vanishes and the
occupation-cancelled (diamagnetic) seagull must be subtracted for a finite answer;
once regularized, the naive current--current object returns the \emph{wrong} tensor
(uniaxial about $(1,-1,0)$, not the optical axis), and the correct assembly is the
orbital-susceptibility (Berry-curvature plus quantum-metric) combination.  This
divergence and its wrong-axis outcome were diagnosed in Paper~VIII
\cite{menendez2026electricblock} and reproduced here in an exploratory band-sum
probe; the geometric route of Sec.~\ref{subsec:blocks} sidesteps them entirely.
Since the geometric and fermion-loop routes
compute the \emph{same} induced action and must agree---and the geometric route
already yields Eq.~\eqref{eq:blocks_explicit}---the dynamical extraction is a matter
of methodological completeness, not of deciding the verdict, and is deferred to
follow-on work.  As with the kinematic doubler accounting
(Remark~\ref{rem:no_new_modes}), nothing in the birefringence verdict depends on it.
\end{remark}


\section{Falsifiability: A Multi-Channel Map}
% falsifiability.tex -- the multi-channel falsifiability map.
%
% Author decision (handoffs 2026-07-18-paper04-falsifiability-do-both-decision +
% -birefringence-null-is-a-passed-falsifier): DO BOTH channels. 2.b (kinematic) is
% concrete here; 2.a (gauge photon) is a contingent OPEN follow-on. Reframe the
% birefringence null as a PASSED falsifiable prediction (not "not a falsifier"),
% with two honesty guardrails. SME placement verified against Kostelecky--Mewes.
%
% Backs the falsifiability audit row (Table~\ref{tab:audit}).

\label{sec:falsifiability}

The program's motivation for this paper is falsifiability, and the optical axis is a
parameter-free target. It is worth being precise about which of the framework's
predictions actually test it, because the most common objection gets the logic
backwards.

\subsection{The birefringence null is a passed test, not an absence of content}
\label{subsec:passed_null}

Discrete-spacetime proposals are often dismissed on the grounds that ``there is no
evidence of birefringence.'' A falsifiable prediction, however, is precisely one that
\emph{forbids} an observable outcome. The framework forbids optical-axis vacuum
birefringence---the doubled transverse root of Theorem~\ref{thm:doubled_root}---and
vacuum birefringence is both directly observable and among the most tightly bounded
quantities in physics. Its non-observation is therefore not silence; it is a risky
prediction the framework has \emph{passed}.

The test moreover \emph{discriminates}. In the Standard-Model Extension (SME)
classification of Lorentz-violating photon
operators~\cite{kosteleckymewes2009,kosteleckyrussell2011}, operators of odd mass
dimension are CPT-odd and \emph{birefringent}; the leading Planck-suppressed such
operator is dimension five, giving a helicity-dependent vacuum birefringence. Generic
discrete, Lorentz-violating, and quantum-gravity models populate this sector, and it
is exactly the astrophysical non-observation of vacuum birefringence that
\emph{excludes} them. This framework instead sits in the CPT-even, non-birefringent
sector: the induced photon action's adjugate closure (Theorem~\ref{thm:adjugate})
cancels the polarization split, placing it with the even-dimension, dispersive-but-not-
birefringent operators (the $c^{(6)}_{(I)jm}$ family). The same datum that refutes the
birefringent rivals therefore \emph{corroborates} this one. The null is not luck; it is
what the framework's operator class requires.

Two guardrails keep this honest. \emph{(a) Retrospectively}, the null is a passed
falsifiable prediction, and a strong one, because the natural alternatives fail the
same test. \emph{(b) Prospectively}, a further non-detection of birefringence no longer
\emph{discriminates} this framework from standard electrodynamics---so the
forward-testable frontier is the dispersion channels below, not birefringence. And the
cancellation is a theorem \emph{conditional} on the constitutive blocks
(Sec.~\ref{sec:gauge_sector}); their derivation as the full dynamical vacuum response
is itself an open item (Sec.~\ref{subsec:gauge_channel}). The statement is thus
corroborated \emph{and} conditional---no more, and no less.

\subsection{The falsifiable content is dispersion, and it splits by field}
\label{subsec:falsifiability_map}

Because birefringence is forbidden, the forward-testable content is the
direction/energy-dependent dispersion, and the $\omega=0$ analysis of
Sec.~\ref{subsec:a1_normalization} forces it to split into distinct channels living in
distinct fields. We map them honestly, per channel:

{\small
\begin{center}
\begin{tabular}{@{}p{3.5cm}p{1.9cm}p{3.6cm}p{3.4cm}@{}}
\hline
\textbf{Channel} & \textbf{Field} & \textbf{Effect} & \textbf{Status} \\
\hline
Optical-axis birefringence & gauge photon & polarization split & \emph{Forbidden}---passed falsifiable prediction (corroborated; conditional on the blocks) \\
Kinematic massive dispersion & spinor (matter) & dim-6 directional dispersion; sidereal & Candidate falsifiable; real, bound $\gg$ Planck (\texttt{PART}) \\
Kinematic massless anisotropy & spinor ($\omega=0$) & magnitude / filtering anisotropy & Real observable, testability open (\texttt{PART}) \\
Gauge-sector photon dispersion & U(1) gauge & common-mode speed anisotropy & Open---derivation + $O_h$ pending; may be null (\texttt{OPEN}) \\
\hline
\end{tabular}
\end{center}
}

\subsection{The kinematic channel (matter sector)}
\label{subsec:kinematic_channel_falsif}

The concrete falsifiable content of this paper is the massive-sector directional
dispersion of Sec.~\ref{sec:kinematic_channel} ($\omega\neq0$). Its low-energy form
$|\tilde H_\text{RGB}|^2 = 1 - \tfrac43|\mathbf{k}_\perp a|^2$ is a
\emph{dimension-six}-scale anisotropy: the fractional effect is $\sim(E a/\hbar c)^2
\sim (E/E_\text{Planck})^2$ if the lattice spacing is near Planckian. Two things must
be said plainly about its reach.

First, the \emph{magnitude} is small and the honest expectation is a bound, not a
detection. Dimension-six suppression is the most suppressed Lorentz-violating case;
even the strongest photon dimension-six time-of-flight bound---LHAASO's observation of
GRB~221009A~\cite{lhaaso2024grb}, which for an isotropic quadratic dispersion implies a
lattice spacing $a \lesssim 3\times10^{-28}$~m ($\sim\!10^{7}\,\ell_\text{Planck}$)---is
weak in absolute terms, and massive particles forgo the cosmological baseline that gives
photons their lever arm. A Planck-scale lattice sits some twenty orders below current
reach; the correctly scoped claim is a limit far above the Planck scale, not a signal.

Second, and more usefully, the distinctive handle is the \emph{signature}, not the
magnitude. A lattice with a fixed spatial axis $\hat{\mathbf{n}}_*$ is anisotropic in a
way no isotropic theory reproduces: as the Earth rotates, the laboratory frame sweeps
through $\hat{\mathbf{n}}_*$, so the predicted effect \emph{modulates at the sidereal
frequency}. In SME language this is the statement that the anisotropic coefficients
($j\neq0$) are non-zero with a fixed, parameter-free angular pattern about
$\hat{\mathbf{n}}_*$~\cite{kosteleckymewes2009}. A sidereal modulation with this
specific pattern is a qualitatively distinctive target that an isotropic
Lorentz-violating model cannot fake, even where the amplitude is small. The natural
observational handles, stated with their caveats rather than as claims, are then: the
sidereal modulation of high-energy propagation; high-energy astrophysical
\emph{neutrinos} as the one massive channel that keeps a cosmological baseline; and the
$\omega=0$ magnitude anisotropy---a direction-dependent \emph{filtering} of amplitudes
(not an energy dispersion), the anisotropic spreading that \texttt{exp\_01} measures.

\subsection{The gauge-sector photon channel (contingent)}
\label{subsec:gauge_channel}

The genuinely \emph{photonic} dispersion---the object an astrophysical dimension-six
photon bound ($c^{(6)}_{(I)jm}$;
\cite{kosteleckymewes2009,lhaaso2024grb}) directly constrains---is \emph{not} delivered
by the $\omega\to0$ limit of the spinor propagator, which does not disperse
(Prop.~\ref{prop:massless_flat}). On this lattice the photon is the induced gauge field,
and its candidate falsifiable dispersion is the gauge-sector \emph{common-mode} speed
anisotropy $\omega^2 = \mathbf{k}^{\!\top}\boldsymbol{\varepsilon}\,\mathbf{k}$
(Sec.~\ref{subsec:common_mode}). We do not assert it as a bound in hand. It is gated on
two derivations still open: (1) the effective-action ($\operatorname{Tr}\ln T$)
legitimization of $\boldsymbol{\varepsilon}$ and $\boldsymbol{\mu}^{-1}$ as the actual
dynamical vacuum response rather than geometric candidates
(Remark~\ref{rem:vertex_not_action}); and (2) the $O_h$-restoration computation
(Remark~\ref{rem:oh_restoration}), which may return that the anisotropy averages to a
null in the continuum. We name it as the open channel we are still deriving, status
\texttt{OPEN}, and pursue it as follow-on work rather than claiming it here.

Taken together, these are the framework's falsifiable content, honestly mapped: a
forbidden birefringence that the data have already corroborated and that excludes the
birefringent rivals; a real but Planck-suppressed massive dispersion whose distinctive
target is a sidereal signature; a real $\omega=0$ filtering anisotropy; and a contingent
photon dispersion still under derivation. The value is not one frontier number but a
structured, per-channel-honest account of exactly what tests the theory.


%% Add more \section{...} \input{sections/...} blocks as needed.
%% Group with comment dividers (e.g., "PART II: DYNAMICS") to keep
%% main.tex skimmable as the paper grows.

\section{Conclusion}
% conclusion.tex

The electric sector of the A=1 induced gauge action is now as concrete as the
magnetic one. Reading the on-site advance $\delta\phi=\omega+V(x)$ as a temporal
link turns the apparent obstruction---``no electric Wilson loop''---into a
temporal plaquette whose holonomy delivers the permittivity
$\varepsilon=P=\sum_a V_a V_a^{\!\top}$, eigenvalues $\{1,4,4\}$, optical axis
$(1,1,-1)$ suppressed: the exact mirror of the magnetic $\{4,4,16\}$. The block is
not only derived but read off the engine's real tick evolution
(\autoref{subsec:electric_block}), so it rests on the coupling the framework
actually implements, not an analogy. The two sectors are bound by an exact
algebraic identity, $\mu^{-1}=\operatorname{adj}(\varepsilon)$
(\autoref{thm:adjugate}), which makes the impedance product isotropic in every
principal direction and thereby forces the photon dispersion to a doubled
transverse root. Gauge-sector vacuum birefringence about $(1,1,-1)$ therefore
\emph{cancels}, and does so independently of the undetermined coupling
normalizations and of the specific lattice geometry. The companion Paper~IV
confirms the verdict from an independent implementation---matching blocks and a
null split to $\sim10^{-15}$---so the two papers stand together on the same result
by two routes.

What changed by the end is that a question Paper~IV could only pose is now
answered: the framework \emph{predicts the absence} of gauge-sector vacuum
birefringence---a positive null, consistent with polarimetry rather than in
tension with it. Two things are deliberately left open. The shared photon speed
carries a factor-$\sim$2 \emph{common-mode} anisotropy about the optical axis (the
axis is the one cube body-diagonal the hop set omits, so the lattice is trigonal
rather than cubic); this is not birefringence, and whether the physical vacuum
$O_h$-restores to an isotropic speed---rather than realizing a single trigonal
domain---is a genuine open question we do not settle here. Separately, the same
induced action admits a momentum-space route (the physical-band $\mathrm{Tr}\ln T$
eigenphase sum): we fix its mode-filling prescription (the physical band) and
diagnose why a naive band sum diverges (an intraband pole on the exactly-flat
optical axis, Ward-cancelled), but the Ward-safe single-object extraction of the
anisotropic \emph{tensor} is left to a possible follow-on; here the tensor
structure is taken from the plaquette geometry and only the scalar magnitude from
the loop. Neither open item bears on the birefringence verdict.

The natural next steps are those two threads---an isotropy-restoration analysis of
the common-mode speed, and the orbital-susceptibility assembly of the dynamical
tensor---together with the one quantity Paper~I left open and this paper inherits,
the absolute one-loop coupling $1/g^2$, which sets the induced-action magnitude but
not the birefringence verdict.


\section*{Acknowledgements}
% acknowledgements.tex

This paper was developed in close collaboration with Anthropic's Claude (Claude
Opus~4.8, working through Claude Code), which carried the digital formalization
throughout: the closed-form kinematic derivations, the symbolic verification of
the adjugate and doubled-root theorems, the companion experiments
(\texttt{exp\_01}--\texttt{exp\_05}) and their numerical checks, and the drafting
and audit-disciplined revision of the manuscript. The physical framework and its
guiding insights---among them the recognition that the whole optical-axis
structure follows from the arbitrary omission of the fourth cube body-diagonal---%
are the author's; Claude's role was to render them into precise notation, working
code, and prose, and to hold every claim to the paper's audit-table standard.

The no-go framing and several of its sharpest corrections were shaped by a
detailed independent review carried out with a large language model, whose
recomputation of the load-bearing results and whose principal criticisms---the
exact matter-sector flat band, the linearly accumulating $\omega=0$ filter, and
the classical antecedents (Cauchy--Binet; impedance-matched media) of the two
central theorems---materially strengthened the paper. Numerical and symbolic work
used NumPy and SymPy together with the \texttt{dcl\_core} lattice engine; the
document was typeset with \LaTeX. Any remaining errors are the author's.


%% ============================================================
%% APPENDICES
%% ============================================================
\appendix

\section{Code and Data}
% code_and_data.tex
\label{sec:code_and_data}

Every load-bearing claim is backed by a runnable script in the companion
repository, and the paper's derived equations are generated from those scripts
(via \texttt{sympy.latex}) into \texttt{paper/sections/generated/} and
\verb|\input| where they appear, so the displayed equations cannot drift from
what was computed.

\paragraph{Derivations (\texttt{src/utilities/}).}
\texttt{electric\_induced\_action.py} is the core: it re-derives the magnetic
anchor $Q_B=\{4,4,16\}$ from first principles, the electric block
$\varepsilon=P=\{1,4,4\}$, the adjugate theorem
$\mu^{-1}=\operatorname{adj}(\varepsilon)$ (for general hop vectors), the doubled-root
dispersion and null split (symbolic, plus a $2\times10^4$-direction numeric sweep),
and the $O_h$ domain average. \texttt{trlnT\_prescription.py} establishes the
$\mathrm{Tr}\ln T$ mode-filling prescription (physical band); \texttt{ward\_safe\_
diagnosis.py} diagnoses the momentum-space band-sum divergence.

\paragraph{Engine cross-check (\texttt{src/experiments/}).}
\texttt{exp\_01\_electric\_permittivity\_extraction.py} reads $\varepsilon=P$ off the
dcl-core engine's real \texttt{HopOperator.step} output (importing
\texttt{dcl\_core.core3d}), the temporal-sector mirror of dcl-core \texttt{exp\_04};
its \texttt{data/} outputs record the extracted tensor and the checks.

\paragraph{Notes (\texttt{notes/}).}
\texttt{electric\_block\_derivation.md} is the narrative derivation;
\texttt{speed\_anisotropy\_and\_isotropy\_restoration.md},
\texttt{trlnT\_prescription.md}, and \texttt{ward\_safe\_vacuum\_polarization.md}
cover the open items. \texttt{paper/sections/audit\_table.tex}
(Table~\ref{tab:audit}) is the authority mapping each claim to its script.


\section{Reproducibility Quickstart}
% reproducibility.tex
\label{sec:reproducibility}

\paragraph{Environment.}
Python (with \texttt{sympy} and \texttt{numpy}); the engine cross-check additionally
imports \texttt{dcl\_core} (dcl-core v0.3.0). No GPU is required---the derivations
are symbolic and the extraction runs on a small lattice.

\paragraph{Verify the headline result.}
From a fresh clone, the central claims are reproduced by
%
\begin{quote}\ttfamily
python -u src/utilities/electric\_induced\_action.py
\end{quote}
%
which prints the magnetic anchor $\{4,4,16\}$, the electric block $\{1,4,4\}$ (axis
suppressed), $\mu^{-1}=\operatorname{adj}(\varepsilon)$, and the null birefringence
split, all \texttt{PASS}. The engine cross-check that $\varepsilon=P$ is what the
tick rule actually produces is
%
\begin{quote}\ttfamily
python -u src/experiments/exp\_01\_electric\_permittivity\_extraction.py
\end{quote}

\paragraph{Build the paper.}
From \texttt{paper/}: \texttt{pdflatex main} $\to$ \texttt{bibtex main} $\to$
\texttt{pdflatex main} (twice). The \verb|\input| equation fragments under
\texttt{sections/generated/} are regenerated by
\texttt{electric\_induced\_action.py}.

\paragraph{Where to look next.}
The open items and the momentum-space route are in
\texttt{src/utilities/trlnT\_prescription.py} and
\texttt{ward\_safe\_diagnosis.py} (with the matching notes); the independent
numerical confirmation of the birefringence verdict is in the companion
Paper~IV~\cite{PaperIV}.


\section{The Ward-Safe Band-Sum}
% band_sum.tex -- Appendix: the Ward-safe band-sum estimate of the
% common-mode anisotropy order (referee M5: bring the load-bearing dynamical
% evidence into the paper). Source: notes/gauge_common_mode_order.md.
\label{sec:band_sum}

The ``strong-provisional'' half of the gauge common-mode verdict
(Sec.~\ref{subsec:gauge_channel}) rests on a dynamical estimate of the
\emph{order} of the anisotropy, independent of the geometric blocks. We record it
here so the claim is auditable from the paper.

\paragraph{The object.} On the validated Bloch transfer operator $T(\mathbf{k})$,
we form the physical-band second-order response to a photon of wavevector
$\mathbf{q}$, with the intraband contribution dropped per the Ward-safe
prescription (the physical band is exactly flat along the optical axis, so the
intraband denominator $\lambda(\mathbf{k})-\lambda(\mathbf{k}\pm\mathbf{q})$
vanishes there and the naive sum is Ward-divergent). Sweeping $|\mathbf{q}|\to0$
gives the axis-versus-perpendicular anisotropy ratio of the resulting tensor:

\begin{center}
\begin{tabular}{@{}lcccc@{}}
\hline
$|\mathbf{q}|$ & $0.30$ & $0.20$ & $0.12$ & $0.07$ \\
\hline
max\,:\,min ratio & $0.60$ & $0.37$ & $0.33$ & $0.32$ \\
\hline
\end{tabular}
\end{center}

\paragraph{What it shows.} The ratio \emph{converges to a constant} $\approx0.32$
(not to $1$) as $|\mathbf{q}|\to0$. A dimension-six--suppressed anisotropy would
have the ratio approach $1$ in that limit; a constant, non-unity limit is the
signature of an $O(1)$, dimension-four anisotropy---the same order the geometric
blocks give directly (a $\mathbf{k}$-independent factor-$\sim2$ speed). Two
independent routes therefore agree that the common-mode anisotropy is not
suppressed.

\paragraph{The caveat (an instructive wrong tensor).} This interband-only object is
\emph{not} the correct induced tensor: it is uniaxial about $(1,-1,0)$, not the
optical axis $(1,1,-1)$, so it is used here only to gauge the \emph{order} of the
anisotropy, not its structure. The correct assembly is the orbital-susceptibility
combination (Berry curvature plus quantum metric), a stated open item
(Remark~\ref{rem:vertex_not_action}). For the correct object to be
dimension-six--suppressed, those extra terms would have to \emph{cancel} the
$O(1)$ anisotropy the naive object exhibits---a non-generic cancellation, and the
available evidence is against it. This is why the dynamical verdict is
\emph{strong-provisional} rather than decisive: the estimate is robust as to order,
but the exact tensor is uncomputed. Computation:
\texttt{notes/gauge\_common\_mode\_order.md} and the associated Ward-safe
$\mathbf{q}$-scan on $T(\mathbf{k})$.


%% ============================================================
\bibliographystyle{unsrt}
\bibliography{paper-bib/references}

\end{document}
